\section{Real Robot Deployment}

The experiments above evaluate the typing step in isolation, on hand-authored problems. To show that this step composes into a working system rather than depending on curated PDDL, we deploy the full pipeline on a physical robot, where the typed problem is constructed end-to-end from raw perception rather than provided a priori.

Figure~\ref{fig:robot_example} shows the real robot domain, deployed on a Stretch mobile manipulator~\cite{stretch}. The domain includes actions such as \texttt{pickup} and \texttt{clean\_whiteboard}, each with typed preconditions that constrain which objects they can be applied to. At runtime, the scene is perceived using Grounding DINO~\cite{gdino} for open-vocabulary object detection and SAM 2~\cite{sam} for segmentation: the detector returns labeled bounding boxes for a set of prompted object classes, and SAM 2 refines these into pixel-level masks used to estimate each object's 3D position from aligned depth. The result is a set of detected objects with semantic class names --- the system knows, for example, that a \texttt{cloth} is present --- but has no knowledge of the affordances required by the planner. Accordingly, all detected objects are assigned the abstract base type \texttt{graspable} in the generated PDDL problem, and the LLM must assign appropriate types before planning can proceed. This is precisely the typing gap our method resolves, now arising organically from perception: the vision stack reports \emph{what} is in the scene, but not \emph{what role} each object can play.

The accompanying video shows the running example of Figure~\ref{fig:robot_example} executed on hardware: with no \texttt{eraser} present, the system types the \texttt{cloth} as a functional substitute and the robot completes the whiteboard-cleaning goal. We report the deployment as an integration demonstration rather than a quantitative benchmark by design. The claim under evaluation --- that the LLM resolves typing gaps correctly --- is measured in the controlled experiments above, where success depends only on the type assignment. Physical execution instead exercises orthogonal subsystems (grasping, navigation, depth estimation) whose failures are independent of typing correctness; folding them into a success metric would confound rather than test the contribution. The robot therefore establishes real-world applicability, while the benchmark establishes correctness.
